\documentclass[11pt]{article}

\usepackage{amsmath,amssymb,amsthm,mathtools}
\usepackage[margin=1in]{geometry}
\usepackage{microtype}

\newtheorem{theorem}{Theorem}[section]
\newtheorem{lemma}[theorem]{Lemma}
\newtheorem{corollary}[theorem]{Corollary}
\newtheorem{remark}[theorem]{Remark}

\newcommand{\dd}{\,\mathrm d}

\title{Cliques with Distributed Pendant Leaves:\\
Symmetric Fractional Degree Bias}
\author{}
\date{}

\begin{document}

\maketitle

\begin{abstract}
Fix $r\geq3$.  Attach $n_i\geq0$ pendant leaves to vertex $i$ of $K_r$
and put $h=n_1+\cdots+n_r$.  We prove that, whenever $1\leq h\leq r$,
every graphon $W$ of edge density
$p\geq1-1/(r-1)$ satisfies
\[
 t(H,W)\geq p^h\prod_{j=1}^{r-1}\bigl(jp-(j-1)\bigr)
 =\Phi_H(p).
\]
The leaf multiplicities may be distributed arbitrarily among the clique
vertices.  The proof averages over clique permutations, uses the
arithmetic--geometric mean inequality to replace the distribution by the
fractional degree weight $d^{h/r}$ at every clique vertex, and changes the
underlying probability measure.  A fractional H\"older inequality and a
monotone rescaling of the accepted clique bound complete the argument.  The
case $(r,h)=(4,2)$ with multiplicities $(1,1,0,0)$ proves the formerly open
NetworkX Atlas graph $134$.
\end{abstract}


%%%%%%%%%%%%%%%%%%%%%%%%%%%%%%%%%%%%%%%%%%%%%%%%%%%%%%%%%%%%%%%%%%%%%%%%%%%%%%%
\section{The graph family and chromatic target}
%%%%%%%%%%%%%%%%%%%%%%%%%%%%%%%%%%%%%%%%%%%%%%%%%%%%%%%%%%%%%%%%%%%%%%%%%%%%%%%

Let $(\Omega,\mu)$ be an arbitrary probability space and let
$W\colon\Omega^2\to[0,1]$ be symmetric and measurable.  Write
\[
 p=\int_{\Omega^2}W(x,y)\dd\mu(x)\dd\mu(y),
 \qquad
 d(x)=\int_\Omega W(x,y)\dd\mu(y).
\]

For integers $r\geq3$ and $n_1,\ldots,n_r\geq0$, let
\[
 H=H_r(n_1,\ldots,n_r)
\]
be obtained from the clique on labelled vertices $1,\ldots,r$ by attaching
$n_i$ private pendant leaves to vertex $i$.  Put
\[
 h=n_1+\cdots+n_r.
\]
The theorem below assumes $1\leq h\leq r$; no restriction is placed on the
individual $n_i$.

Since a leaf has $x-1$ colour choices after its neighbour has been coloured,
\begin{equation}
 \chi_H(x)=(x)_r(x-1)^h,
 \qquad
 (x)_r=x(x-1)\cdots(x-r+1).
 \label{eq:chromatic}
\end{equation}
The clique is a subgraph and the leaves do not increase chromatic number, so
$\chi(H)=r$.  Define the usual clique target
\[
 A_r(p):=\prod_{j=1}^{r-1}\bigl(jp-(j-1)\bigr).
\]
Substituting $x=(1-p)^{-1}$ in \eqref{eq:chromatic} gives
\begin{equation}
 \Phi_H(p)=p^hA_r(p).
 \label{eq:target}
\end{equation}
Both sides are polynomials in $p$, so this identity includes the continued
value $1$ at $p=1$.  At the chromatic threshold
\[
 a_r:=1-\frac1{r-1}=\frac{r-2}{r-1},
\]
the final factor of $A_r$ vanishes and hence $\Phi_H(a_r)=0$.

We use exactly one previously accepted graphon theorem: for every graphon
$V$ of edge density $s\geq a_r$,
\begin{equation}
 t(K_r,V)\geq A_r(s).
 \label{eq:clique-input}
\end{equation}
This is the complete-clique instance of
\texttt{notes/pure\_chordal.tex} (equivalently, the accepted Moon--Moser
clique bound).  It applies to graphons on arbitrary probability spaces.


%%%%%%%%%%%%%%%%%%%%%%%%%%%%%%%%%%%%%%%%%%%%%%%%%%%%%%%%%%%%%%%%%%%%%%%%%%%%%%%
\section{Fractional degree bias}
%%%%%%%%%%%%%%%%%%%%%%%%%%%%%%%%%%%%%%%%%%%%%%%%%%%%%%%%%%%%%%%%%%%%%%%%%%%%%%%

We record the weighted-edge estimate needed after changing measure.

\begin{lemma}[Fractionally degree-weighted edge]
\label{lem:fractional-edge}
Let $0<\alpha\leq1$ and set
\[
 N_\alpha
 =\int_{\Omega^2}W(x,y)d(x)^\alpha d(y)^\alpha
   \dd\mu(x)\dd\mu(y).
\]
Then
\[
 N_\alpha\geq p^{1+2\alpha}.
\]
\end{lemma}

\begin{proof}
On $\{d=0\}$, define $W(x,y)/d(x)$ to be zero.  Fubini's theorem implies
that $W(x,y)=0$ for almost every $y$ there.  Put
\[
 \theta=\frac1{1+2\alpha},
 \qquad
 \eta=\frac{\alpha}{1+2\alpha}.
\]
Thus $\theta+2\eta=1$ and, almost everywhere on the support of $W$,
\[
 W(x,y)
 =\bigl(W(x,y)d(x)^\alpha d(y)^\alpha\bigr)^\theta
  \left(\frac{W(x,y)}{d(x)}\right)^\eta
  \left(\frac{W(x,y)}{d(y)}\right)^\eta.
\]
Weighted H\"older gives
\[
 p\leq N_\alpha^\theta
 \left(\int\frac{W(x,y)}{d(x)}\dd\mu^2\right)^\eta
 \left(\int\frac{W(x,y)}{d(y)}\dd\mu^2\right)^\eta.
\]
Each quotient integral is the measure of a subset of $\{d>0\}$ and is at
most one.  Hence $p\leq N_\alpha^\theta$, which is the assertion.
\end{proof}

The next scalar lemma is the reason that the upper bound $h\leq r$ is the
natural boundary of this argument.

\begin{lemma}[Clique rescaling]
\label{lem:scalar}
Let $r\geq3$, let $1\leq h\leq r$, and let $p\in(a_r,1]$.  Suppose
$0<z\leq p^h$ and put
\[
 s_0=p\left(\frac{p^h}{z}\right)^{2/r}.
\]
If $s_0\leq1$, then
\[
 zA_r(s_0)\geq p^hA_r(p).
\]
\end{lemma}

\begin{proof}
Every factor $js-(j-1)$ is positive for $s>a_r$.  Put
\[
 \Psi_r(s)=\frac{A_r(s)}{s^{r/2}}.
\]
For $s\in(a_r,1]$, logarithmic differentiation gives
\begin{align*}
 \frac{\Psi_r'(s)}{\Psi_r(s)}
 &=\sum_{j=1}^{r-1}\frac{j}{js-(j-1)}-\frac{r}{2s}\\
 &=\frac{r-2}{2s}
   +\sum_{j=2}^{r-1}
      \frac{j-1}{s\bigl(js-(j-1)\bigr)}>0.
\end{align*}
Thus $\Psi_r$ is strictly increasing.  The definition of $s_0$ is
equivalent to
\[
 z=p^h\left(\frac p{s_0}\right)^{r/2}.
\]
Since $s_0\geq p$, it follows that
\[
 zA_r(s_0)
 =p^{h+r/2}\Psi_r(s_0)
 \geq p^{h+r/2}\Psi_r(p)
 =p^hA_r(p).
\]
\end{proof}

Also, $A_r$ is nondecreasing on $[a_r,1]$: every one of its factors is
nonnegative and nondecreasing there.


%%%%%%%%%%%%%%%%%%%%%%%%%%%%%%%%%%%%%%%%%%%%%%%%%%%%%%%%%%%%%%%%%%%%%%%%%%%%%%%
\section{The distributed-leaf theorem}
%%%%%%%%%%%%%%%%%%%%%%%%%%%%%%%%%%%%%%%%%%%%%%%%%%%%%%%%%%%%%%%%%%%%%%%%%%%%%%%

\begin{theorem}[Leaves distributed on a clique]
\label{thm:main}
Let $r\geq3$ and let $n_1,\ldots,n_r\geq0$ have total
$1\leq h\leq r$.  For every graphon $W$ of edge density
\[
 p\in[a_r,1],
\]
one has
\[
 \boxed{
 t\bigl(H_r(n_1,\ldots,n_r),W\bigr)
 \geq p^hA_r(p)
 =\Phi_{H_r(n_1,\ldots,n_r)}(p).
 }
\]
\end{theorem}

\begin{proof}
The claim is immediate at $p=a_r$, because the target is zero and every
homomorphism density is nonnegative.  At $p=1$, one has $W=1$ almost
everywhere, and both sides equal one.  Hence assume $p\in(a_r,1)$.

For clique variables $x_1,\ldots,x_r$, write
\[
 F(\boldsymbol x)=\prod_{1\leq i<j\leq r}W(x_i,x_j).
\]
After integrating every leaf, the homomorphism density is
\[
 t(H,W)=\int_{\Omega^r}F(\boldsymbol x)
              \prod_{i=1}^r d(x_i)^{n_i}\dd\mu^r.
\]
Permuting the clique variables shows that the same integral is obtained
after any permutation of the exponent vector $(n_1,\ldots,n_r)$.  Average
over all $r!$ permutations.  The arithmetic--geometric mean inequality,
applied pointwise to the resulting $r!$ monomials, gives
\begin{equation}
 t(H,W)
 \geq\int_{\Omega^r}F(\boldsymbol x)
       \prod_{i=1}^r d(x_i)^{h/r}\dd\mu^r.
 \label{eq:symmetrization}
\end{equation}
Indeed, each $n_j$ appears in each coordinate exactly $(r-1)!$ times, so
the geometric-mean exponent of every $d(x_i)$ is
$(r-1)!h/r!=h/r$.

Set
\[
 \alpha=\frac hr\in(0,1],
 \qquad
 M=\int_\Omega d(x)^\alpha\dd\mu(x),
 \qquad
 z=M^r.
\]
Since $p>0$, one has $M>0$.  Define the probability measure
\[
 \dd\nu(x)=\frac{d(x)^\alpha}{M}\dd\mu(x).
\]
The right side of \eqref{eq:symmetrization} is
\begin{equation}
 z\,t(K_r,W;\nu).
 \label{eq:measure-change}
\end{equation}
The edge density of $W$ on $(\Omega,\nu)$ is
\[
 s=\frac{N_\alpha}{M^2}.
\]

Concavity of $u\mapsto u^\alpha$ gives
\[
 M\leq p^\alpha,
 \qquad
 0<z\leq p^h.
\]
By Lemma~\ref{lem:fractional-edge},
\[
 s\geq\frac{p^{1+2h/r}}{M^2}
 =p\left(\frac{p^h}{z}\right)^{2/r}
 =:s_0.
\]
Because $W\leq1$ and $\nu$ is a probability measure, $s\leq1$, so
$s_0\leq1$.  Also $s\geq s_0\geq p>a_r$.

Apply the accepted clique bound \eqref{eq:clique-input} on the biased
probability space, monotonicity of $A_r$, and Lemma~\ref{lem:scalar}:
\begin{align*}
 t(H,W)
 &\geq z\,t(K_r,W;\nu)\\
 &\geq zA_r(s)
 \geq zA_r(s_0)
 \geq p^hA_r(p).
\end{align*}
This is exactly \eqref{eq:target}.
\end{proof}

Every balanced complete $k$-partite graphon with $k\geq r-1$ gives the
expected equality.  At $k=r-1$ both sides vanish.  For $k\geq r$, the degree
is constant, equality holds in the symmetrization and fractional H\"older
steps, the biased measure equals the original measure, and the accepted
clique bound is sharp.


%%%%%%%%%%%%%%%%%%%%%%%%%%%%%%%%%%%%%%%%%%%%%%%%%%%%%%%%%%%%%%%%%%%%%%%%%%%%%%%
\section{Atlas 134, audits, and exact limitations}
%%%%%%%%%%%%%%%%%%%%%%%%%%%%%%%%%%%%%%%%%%%%%%%%%%%%%%%%%%%%%%%%%%%%%%%%%%%%%%%

\begin{corollary}[Atlas 134]
NetworkX Atlas graph $134$, with graph6 code \texttt{E\char126\_O}, is
positive.  It has order six, size eight, chromatic number four, and
\[
 \chi_{F_{134}}(x)=x(x-1)^3(x-2)(x-3).
\]
For every graphon of edge density $p\in[2/3,1]$,
\[
 t(F_{134},W)
 \geq p^3(2p-1)(3p-2)
 =\Phi_{F_{134}}(p).
\]
\end{corollary}

\begin{proof}
Its NetworkX edge list is
\[
 01,\ 02,\ 03,\ 04,\ 12,\ 13,\ 23,\ 35.
\]
Vertices $0,1,2,3$ induce $K_4$, while $4$ and $5$ are leaves attached to
the distinct clique vertices $0$ and $3$.  Thus the graph is
$H_4(1,0,0,1)$, and Theorem~\ref{thm:main} applies with $h=2$.
\end{proof}

All integrations above concern bounded nonnegative measurable functions, so
Tonelli--Fubini applies without an approximation or finite-step reduction.
The quotient $W/d$ is defined on $\{d=0\}$ and the almost-everywhere
vanishing needed there is stated explicitly.  The biased measure is formed
only after proving $M>0$.  The proof uses ordinary homomorphism densities,
not injective copies, and assumes neither regularity nor a minimum-degree
condition.  The threshold and $p=1$ endpoints are handled before any use of
positive clique factors or logarithmic differentiation.

The theorem's exact new boundary is $h\leq r$.  This is what ensures that
$d^{h/r}$ is a concave weight and that Lemma~\ref{lem:fractional-edge}
applies with exponent at most one.  The result does not cover more than $r$
leaves distributed over several roots.  It also requires the non-leaf core
to be a complete clique and every attachment to be a pendant leaf; it is not
a general leaf-attachment or mixed-chordal closure theorem.

The finite chromatic identities, scalar derivative, H\"older exponents, and
Atlas identification are reconstructed exactly by
\begin{verbatim}
python codes/verify_clique_distributed_leaves.py
\end{verbatim}
For formal verification, the reusable analytic pieces are the permutation
symmetrization \eqref{eq:symmetrization}, the fractional edge lemma, and the
scalar monotonicity of $A_r(s)/s^{r/2}$.  The graph-specific module awaiting
inspection is
\texttt{lean/Taeyoung/Examples/Graph134.lean}.


%%%%%%%%%%%%%%%%%%%%%%%%%%%%%%%%%%%%%%%%%%%%%%%%%%%%%%%%%%%%%%%%%%%%%%%%%%%%%%%
\section{What the Lean formalization actually proves}
%%%%%%%%%%%%%%%%%%%%%%%%%%%%%%%%%%%%%%%%%%%%%%%%%%%%%%%%%%%%%%%%%%%%%%%%%%%%%%%

The machine-checked development is
\texttt{lean/Taeyoung/Methods/CliqueDist/\{Bias,Rows\}.lean}, ending at
\texttt{satisfiesLowerBound\_134}.  It proves the case
$(r,h)=(4,2)$ with multiplicities $(1,1,0,0)$, and only that case.  Three
things are done differently.

\paragraph{Scope: $r=4$, $h=2$, and $\alpha=1/2$.}
Theorem~\ref{thm:main} is proved here for every $r\geq3$ and every
$1\leq h\leq r$.  The formalization fixes $r=4$, $h=2$.  The reason is the
exponent: $\alpha=h/r$ is $1/2$ there, so the biased density is $\sqrt d$ and
every step stays inside `Real.sqrt`; a general $\alpha$ would need
`Real.rpow` throughout and a fractional H\"older with $r$ factors.  No other
member of the family is in the six-vertex catalogue.

\paragraph{The permutation average becomes a single exchange.}
Equation~\eqref{eq:symmetrization} averages the density over all $r!$
permutations of the exponent vector and applies AM--GM to the resulting $r!$
monomials.  At $(4,2)$ the formalization uses \emph{one} exchange instead.
Writing $H_{ij}$ for $K_4$ with its two leaves at clique vertices $i$ and $j$,
the graphs $H_{01}$ and $H_{23}$ are isomorphic, so
\[
 2\,t(H,W)=\int F\bigl(d_0d_1+d_2d_3\bigr)
 \geq\int F\cdot2\sqrt{d_0d_1d_2d_3}
 =2\int F\prod_{i}\sqrt{d_i},
\]
where the middle step is the two-term AM--GM $a+b\geq2\sqrt{ab}$ applied to the
single perfect matching $\{01\},\{23\}$.  That produces exactly the geometric
weight $\prod d_i^{h/r}$ of \eqref{eq:symmetrization}, from one graph
isomorphism (\texttt{CliqueDist.isoExchange}, checked by \texttt{decide}) rather
than an $r!$-fold average.  The general case still needs the average; the
matching shortcut is available whenever $h=r/2$ and the leaves sit on a
perfect matching of the clique.

\paragraph{Lemma~\ref{lem:scalar} becomes a factorization.}
This note proves the rescaling by showing $\Psi_r(s)=A_r(s)/s^{r/2}$ strictly
increasing, through logarithmic differentiation and the identity
\[
 \frac{\Psi_r'}{\Psi_r}
 =\frac{r-2}{2s}+\sum_{j=2}^{r-1}\frac{j-1}{s(js-(j-1))}>0 .
\]
At $r=4$ the formalization uses no derivative.  With
$m=M^2$ and $s$ the biased edge density, the two inputs are $m\leq p$ and
$ms\geq p^2$; then $m^2s^2\geq p^4$ and
\begin{equation}
 p(2s-1)(3s-2)-s(2p-1)(3p-2)=(s-p)(6ps-2),
 \label{eq:lean-rescale}
\end{equation}
both of whose factors are nonnegative because $s\geq p\geq2/3$.  Multiplying
out gives $m^2A_4(s)\geq p^2A_4(p)$ directly.  Identity
\eqref{eq:lean-rescale} is the $r=4$ shadow of the monotonicity of $\Psi_4$,
and it is a `ring` step.

\paragraph{Everything else transcribes directly.}
Lemma~\ref{lem:fractional-edge} at $\alpha=1/2$ is
\texttt{GeometricMean.sq\_le\_integral\_geoIntegrand}, already proved for
\texttt{notes/two\_root\_triangle\_leaf\_cones.tex}: it is precisely
$N_{1/2}=\int\!\!\int W\sqrt{d(x)d(y)}\geq p^2$, and the three-factor H\"older
of this note's proof is the same one used there.  The concavity bound
$M\leq p^{\alpha}$ becomes $M^2\leq p$, one application of the unweighted
Cauchy--Schwarz \texttt{BookTail.sq\_integral\_mul\_le} against the constant
$1$ --- no concavity is needed.  The change of measure
\eqref{eq:measure-change} is \texttt{CliqueDist.integral\_sqrtDegree\_prod},
built on \texttt{Foundation/TiltTransfer.lean} exactly as
\texttt{Methods/DegreeBias.lean} is at $\alpha=1$.  The clique input
\eqref{eq:clique-input} is the accepted pure-chordal bound
\texttt{cliqueDensity\_ge\_cliquePoly} at $r=4$.

\end{document}
